% ----------------------------------------------------------
% Conclusão
% ----------------------------------------------------------
\chapter[Conclusão]{Conclusão}
%\addcontentsline{toc}{chapter}{Conclusão}
% ---

Este trabalho teve o objetivo de desenvolver uma ferramenta para auxiliar os alunos dos cursos da área de computação no aprendizado de linguagens formais e máquinas de estados, focando principalmente em autômatos de pilha não determinísticos (APNs), dado que esses conhecimentos são fundamentais, porém complexos e abstratos.

Entre as principais contribuições do projeto desenvolvido estão: a criação de APNs; a visualização dos múltiplos caminhos de execução de um APN passo a passo; a criação de correção de exercícios; tudo em uma plataforma completamente online, sem necessidade de instalação no computador.

Apesar disso, o sistema possui algumas limitações, sendo a execução de autômatos que podem criar um número exponencial de múltiplos caminhos quando computados com palavras grandes demais. Essa limitação varia muito de acordo com o hardware do usuário, já que é executado completamente em frontend, porém isso teoricamente não chega a ser um grande problema, já que o sistema foi desenvolvido pensando em executar autômatos voltados a exemplos didáticos e não em simulações de grande porte para pesquisas.

Dessa forma, conclui-se que a ferramenta desenvolvida demonstra potencial de auxiliar os alunos no melhor entendimento acerca de autômatos de pilha não determinísticos, porém não foi realizada pesquisa com os alunos para determinar sua eficiência ou efetividade, sendo esse um dos principais tópicos a serem abordados em trabalhos futuros, juntamente com a implementação de outras máquinas de estados além dos APNs e, talvez, o desenvolvimento de um backend para o sistema, que poderia melhorar a performance do sistema e permitir a execução de autômatos maiores.
